
Foundation models require post-hoc optimizations along multiple efficiency dimensions: low-rank adaptation (LoRA) for parameter efficiency, KV cache compression for memory efficiency, and attention sparsification for computational efficiency. Recent work demonstrates that each dimension can be optimized independently—ARD-LoRA achieves 99.3\% of full fine-tuning performance with 0.32\% trainable parameters, KV-CAT enables training for cache compressibility, and Mamba-2 achieves 2-8× speedup via sub-quadratic architectures.

We investigated whether these efficiency properties could be unified through a single structural constraint: residual-corrected Jacobian stable rank during pretraining. The theoretical basis is that Jacobian stable rank (sr_ℓ^res = ||J̃_ℓ||_F^2 / ||J̃_ℓ||_2^2, where J̃_ℓ = J_ℓ - I) should couple to adaptation efficiency via Fisher information (F ≈ J^T J), memory efficiency via activation covariance (Σ_{ℓ+1} ≈ J_ℓ Σ_ℓ J_ℓ^T), and computational efficiency via attention entropy.

The implementation attempted gradient-based regularization by adding λ · mean(sr_ℓ^res) to the causal language modeling loss. Stable rank was estimated via Hutchinson trace (10 Rademacher probes) for ||J̃_ℓ||_F^2 and power iteration (5 iterations) for ||J̃_ℓ||_2. All computations were differentiated through PyTorch autodiff to enable gradient-based control.

Results show comprehensive failure. The regularized model produced zero stable rank measurements across all layers at all checkpoints. Training perplexity increased from baseline 59.3 to 45,792 (+77,065\%). Regularization losses reached -17.5 billion despite adaptive lambda tuning. The baseline model converged successfully under identical conditions, confirming the failure is specific to the regularization mechanism.

This paper documents the failure modes, provides root cause analysis, and discusses implications for gradient-based spectral control in neural networks.
